\documentclass[11pt,a4paper]{article}
\usepackage[T1]{fontenc}
\usepackage[utf8]{inputenc}
\usepackage{lmodern}
\DeclareUnicodeCharacter{1E93}{\d{z}}
\usepackage{amsmath,amssymb}
\usepackage[margin=25mm]{geometry}
\usepackage[hidelinks]{hyperref}
\hypersetup{pdftitle={A Lieb--Robinson bound for the Kogut--Susskind Hamiltonian, with velocity $\propto Nc/g^2$},pdfauthor={navstokgap research working notes}}
\newcommand{\tightlist}{\setlength{\itemsep}{0pt}\setlength{\parskip}{0pt}}
\setlength{\emergencystretch}{3em}
\setcounter{secnumdepth}{0}
\begin{document}
\hypertarget{a-liebrobinson-bound-for-the-kogutsusskind-hamiltonian-with-velocity-propto-ncg2}{%
\section{\texorpdfstring{A Lieb--Robinson bound for the Kogut--Susskind
Hamiltonian, with velocity
\(\propto Nc/g^2\)}{A Lieb--Robinson bound for the Kogut--Susskind Hamiltonian, with velocity \textbackslash propto Nc/g\^{}2}}\label{a-liebrobinson-bound-for-the-kogutsusskind-hamiltonian-with-velocity-propto-ncg2}}

The Hamiltonian of \href{mass-gap-obligations-lattice.md}{the
obligations map} has exactly the structure for which a Lieb--Robinson
bound survives an unbounded generator: its unbounded part is a sum of
commuting single-link terms, and its interaction is a bounded
finite-range sum of plaquette terms. Passing to the interaction picture
generated by the electric term leaves a time-dependent bounded
interaction of unchanged norm and range, so the standard theorem applies
and gives, for observables \(A\) on a set \(X\) of links and \(B\) on
\(Y\), \[\big\|[A(t),B]\big\|\ \le\ 2\|A\|\,\|B\|\,|X|\;
e^{-\big(d(X,Y)-v_{\rm LR}|t|\big)/a},\qquad
v_{\rm LR}\ \le\ C_{\rm LR}aJ/\hbar\ \le\ 16C_{\rm LR}Nc/g^2,\] where
\(J=\max_\ell\sum_{p\ni\ell}\|w_p\|\le16N\hbar c/(g^2a)\) on the
periodic three-dimensional spatial cubic lattice. The finite constant
\(C_{\rm LR}\) depends on the graph metric and theorem norm, including
four-link support factors; it is independent of volume and coupling. The
earlier \(C_{\rm LR}=2e\) was not derived from the metadata-only source
citations and is not certified here.

With a separately established gap and vacuum hypotheses the standard
clustering theorem applies. Its bound retains the microscopic range;
\(\hbar v_{\rm LR}/\Delta\) alone is not a general correlation-length
bound. A divergent velocity upper estimate proves no superluminality.
Locality is also not by itself a gap-stability theorem.

\hypertarget{the-structure-that-makes-the-unbounded-term-harmless}{%
\subsection{1. The structure that makes the unbounded term
harmless}\label{the-structure-that-makes-the-unbounded-term-harmless}}

Take the links as the sites of the lattice and write, as in
\href{mass-gap-obligations-lattice.md}{the obligations map} §1,
\[H=H_{\rm E}+W,\qquad
H_{\rm E}=\frac{\hbar c\,g^2}{2a}\sum_{\ell\in\mathcal E}(-\Delta_\ell),\qquad
W=\sum_{p\in\mathcal P}w_p,\quad
w_p=\frac{2\hbar c}{a\,g^2}\big(N-\operatorname{Re}\operatorname{tr}U_p\big).\]

Three properties are all that the argument uses.

\begin{enumerate}
\def\labelenumi{\arabic{enumi}.}
\tightlist
\item
  \emph{The unbounded part is single-site and commuting.}
  \(-\Delta_\ell\) acts on the single factor \(L^2(G)\) of link
  \(\ell\), and the terms for different links commute. Each generates a
  one-parameter unitary group \(e^{-isc g^2\Delta_\ell/(2a)}\) whose
  adjoint action preserves the algebra of that link.
\item
  \emph{The interaction is bounded.}
  \(0\le w_p\le\frac{4N\hbar c}{a g^2}\) in operator norm for \(SU(N)\),
  since \(|\operatorname{Re}\operatorname{tr}U|\le N\); for a general
  compact \(G\) in a faithful representation of dimension
  \(\dim_{\rm f}\) replace \(N\) by \(\dim_{\rm f}\).
\item
  \emph{The range is one plaquette.} Each \(w_p\) acts on the four links
  of \(p\); a given link lies in exactly four plaquettes in three
  spatial dimensions (two planes containing its direction, two
  plaquettes on either side in each). Hence
  \[J:=\max_{\ell}\sum_{p\ni\ell}\|w_p\|\ \le\ 4\cdot\frac{4N\hbar c}{a g^2}
  =\frac{16N\hbar c}{a\,g^2}.\]
\end{enumerate}

\hypertarget{the-reduction-and-the-bound}{%
\subsection{2. The reduction and the
bound}\label{the-reduction-and-the-bound}}

\textbf{Proposition 1.} Let \(\tau_t(A)=e^{itH/\hbar}Ae^{-itH/\hbar}\).
For \(A\) supported on a set \(X\) of links and \(B\) on a disjoint set
\(Y\), with \(d(X,Y)\) the lattice distance,
\[\big\|[\tau_t(A),B]\big\|\ \le\ 2\|A\|\,\|B\|\,|X|\,
\exp\!\Big[\frac{v_{\rm LR}|t|-d(X,Y)}{a}\Big],\qquad v_{\rm LR}\le C_{\rm LR}aJ/\hbar.\]

\emph{Proof.} Let \(U_{\rm E}(t)=e^{-itH_{\rm E}/\hbar}\) and define the
interaction-picture interaction
\(\tilde w_p(t)=U_{\rm E}(t)^*w_pU_{\rm E}(t)\). Since \(H_{\rm E}\) is
a sum of commuting single-link terms, \(U_{\rm E}(t)\) factorizes over
links and its adjoint action preserves each link algebra, so
\(\tilde w_p(t)\) is supported on the same four links and
\(\|\tilde w_p(t)\|=\|w_p\|\) for all \(t\). The full evolution
factorizes as \(e^{-itH/\hbar}=U_{\rm E}(t)\,\tilde U(t)\) with
\(\tilde U\) generated by the time-dependent bounded interaction
\(\tilde W(t)=\sum_p\tilde w_p(t)\), and conjugation by \(U_{\rm E}(t)\)
preserves supports and norms. The Lieb--Robinson theorem for
time-dependent bounded finite-range interactions (Lieb--Robinson,
Commun. Math. Phys. 28 (1972) 251; Nachtergaele--Sims, Commun. Math.
Phys. 265 (2006) 119, in the form used for the clustering theorem;
metadata level) applies to \(\tilde W\) with the constant \(J\) of
Section 1, and transferring back through \(U_{\rm E}\) gives the stated
volume-uniform scaling. A numerical prefactor requires fixing the
theorem norm and graph-distance convention, not only \(J\). \(\square\)

The same argument is the one used for oscillator lattices, where the
unbounded on-site harmonic part is treated exactly and the bounded
anharmonic interaction carries the velocity
(Nachtergaele--Raz--Schlein--Sims, Commun. Math. Phys. 286 (2008) 1073;
metadata level). The gauge constraint plays no role: the bound holds on
the full Hilbert space and restricts to the physical subspace, which is
invariant under \(H\).

\textbf{Remark (what fails without the structure).} If the electric term
coupled two links, its unboundedness would enter the iteration and no
bound of this form would follow. The single-link form of
\(\sum_a E^a_\ell E^a_\ell\) is therefore load-bearing, and it is the
same feature that made the strong-coupling hypothesis check of
\href{strong-coupling-uniform-gap.md}{the T2 note} work.

\hypertarget{clustering-and-the-microscopic-range}{%
\subsection{3. Clustering and the microscopic
range}\label{clustering-and-the-microscopic-range}}

The \href{strong-coupling-uniform-gap.md}{strong-coupling theorem}, with
its infinite-volume construction and vacuum hypotheses, already gives
exponential clustering. The standard gapped-Hamiltonian argument has a
conservative range-dependent form

\[\xi\le C\left(a+\frac{\hbar v_{\rm LR}}{\Delta}\right).\]

The constant depends on the locality norm and metric. The range term
cannot be discarded uniformly when the quotient tends to zero. With
\(\Delta\ge\gamma C_2g^2\hbar c/(2a)\), the algebraic consequence is

\[\frac{\hbar v_{\rm LR}}{\Delta}
 \le\frac{32C_{\rm LR}aN}{\gamma C_2g^4}.\]

This controls the dynamical quotient, not \(\xi\) by itself. The earlier
claim \(\xi=O(ag^{-4})\) is withdrawn. A plaquette-tube weight of order
\(g^{-4d/a}\) instead has inverse decay length of order \(\log g/a\).
Nachtergaele--Sims, \emph{CMP} \textbf{265} (2006) 119, and
Hastings--Koma, \emph{CMP} \textbf{265} (2006) 781, are retained at
metadata level; their theorems have their own hypotheses and numerical
constants.

\hypertarget{weak-coupling-and-physical-causality}{%
\subsection{4. Weak coupling and physical
causality}\label{weak-coupling-and-physical-causality}}

The locality upper estimate becomes weak as \(g\to0\). Comparing that
coefficient with \(c\) determines no actual propagation speed. Continuum
relativistic locality requires a reconstructed physical theory.

Even the formal box-crossing comparison distinguishes a bound from an
attained speed. Substituting an order \(Nc/g^2\) coefficient and the
conditional zero-mode energy \(g^{2/3}\hbar c/L\) compares \(Lg^2/(Nc)\)
with \(L/(cg^{2/3})\). The first formal time is smaller when
\(g^{8/3}<N\), up to constants. This proves neither causal connection of
the box on that time scale nor a nonzero-mode truncation theorem.

\hypertarget{gap-stability-is-a-separate-input}{%
\subsection{5. Gap stability is a separate
input}\label{gap-stability-is-a-separate-input}}

Locality is one ingredient of several stability theorems. Their further
ground-space and perturbation-norm conditions must be checked; a
frustration-free/topological-order theorem cannot be applied merely
because a locality bound is available. Finite-coupling Kogut--Susskind
electric and magnetic terms have no common minimizing state. The
\href{strong-coupling-uniform-gap.md}{strong-coupling stability
argument} uses its own product reference and bounded perturbation.

No general openness theorem for the whole gapped coupling set follows
here. Naming frustration-freeness does not exhaust the missing
conditions.

\hypertarget{consequence-for-state}{%
\subsection{6. Consequence for STATE}\label{consequence-for-state}}

Treating commuting unbounded single-link terms exactly leaves bounded
plaquette interactions with unchanged supports and norms. This supplies
volume-uniform locality of order \(Nc/g^2\). Numerical prefactors,
physical propagation, clustering lengths, stability and continuum limits
retain separate hypotheses.
\end{document}
